\chapter{Introduction} \label{Chap1}\label{chap:introduction}
\pagenumbering{arabic}



\section{Motivation and problem statement}

Deep generative models conditioned on widely available satellite imagery offer
a promising route to detailed urban geometry.  Their outputs are not,
however, immediately suitable for an editable city-modelling environment.
Inference is performed over bounded image windows; neighbouring predictions
may disagree in position, height, and topology; and a mesh file does not by
itself identify its geographic extent, building ownership, coordinate frame,
or processing history.  These issues become particularly important when the
output must support procedural editing or interactive use rather than a single
rendered view.  Figure~\ref{fig:sat3dgen-tile-example} illustrates the form of
one retained local output.

\begin{figure}[H]
\centering
\includegraphics[width=0.98\textwidth]{sat3dgen_tile_example.pdf}
\caption[Representative Sat3DGen mesh from one London inference tile.]{Representative
Sat3DGen output for a London inference tile centred at $51.510303^{\circ}$~N,
$0.133978^{\circ}$~W.  Panel (a) renders the stored predicted vertex colours;
panel (b) renders the same vertices and triangles with a uniform material.  The
paired views separate predicted appearance from the generated surface in the
model's normalised local frame.}
\label{fig:sat3dgen-tile-example}
\end{figure}

Procedural city systems address a complementary part of the problem.  Building
footprints, semantic parts, and explicit rules provide stable identity and
authoring control, but do not recover the observed appearance and geometric
detail of a real location automatically.  This dissertation studies the
interface between these representations.  It asks how learned urban geometry
can be combined with open geospatial context and delivered as building-level
assets to ChordAtlas, an existing procedural urban-modelling environment
\cite{kellyChordAtlas2024}.

A case study is centred on London and uses Sat3DGen as the
satellite-conditioned generator \cite{qian2026sat3dgen}.  The study extends its large-image workflow with continuous density placement, smooth
overlap weighting, and vertex-colour recovery before connecting the resulting
surface to downstream geospatial processing.  OpenStreetMap (OSM)
provides footprints and contextual semantics, while an Environment Agency
digital surface model (DSM) supplies an external vertical reference
\cite{openstreetmapData,environmentAgencyDSM2022}.  The resulting problem is
both geometric and operational: tiles must be placed and combined coherently,
and every published asset must retain enough spatial and processing metadata
to be interpreted by the downstream system.

\section{Research aim and question}

The aim is to develop and evaluate a hybrid machine-learning/procedural
workflow that reconstructs a contiguous region from overlapping
satellite-conditioned predictions and converts the resulting geometry into
coordinate-explicit, building-associated assets for ChordAtlas.  The study
addresses the following research question:

\begin{quote}
\textbf{To what extent can overlapping satellite-conditioned predictions be
combined into a coloured regional mesh and integrated with geospatial context
in ChordAtlas through a manifest-driven, coordinate-explicit workflow?}
\end{quote}

The question treats integration as more than file conversion.  A useful result
must establish a consistent spatial frame, associate geometry with geographic
features, preserve appearance information, expose processing state, and
support repeatable transfer between Python mesh processing and the Java-based
ChordAtlas environment.

\section{Design rationale and evaluated scope}

The broader project theme concerned geographically steered world generation
for interactive and collaborative environments.  Sat3DGen was selected because
it produces explicit coloured geometry from satellite imagery, making it a
more direct input to an editing system than an image-only generator.  Initial
feasibility work established local model inference and explored Unity/Ubiq
integration.  A London cross-view collection was then prepared to investigate
location-specific data requirements.

Resource analysis showed that reproducing the released Sat3DGen training
regime would dominate the available computational budget.  The empirical study
therefore retains the released model and evaluates two complementary forms of
regional assembly.  The first adapts the released large-image wrapper so that
rectangular mosaics are fused smoothly in density space and colour is recovered
on the global surface.  The second places independently extracted meshes in a
shared geographic frame and evaluates stitching, height reconciliation, and
building association.  Both paths subsequently enter the same publication
architecture.

OSM and DSM data serve different roles in this design.  Footprints provide
building identity and a basis for per-building extraction; the DSM provides a
metric surface-height reference for vertical processing.  ChordAtlas provides
the procedural destination, including its tiled MiniMesh representation and
building-generation workflow.  BigSUR and FrankenGAN inform the wider goal of
structured fa\c{c}ade modelling, while CityEngine provides a procedural
workflow comparator and Ubiq motivates future collaborative deployment
\cite{kelly2017bigsur,kelly2018frankengan,esriCityEngineIntro2026,
friston2021ubiq}.  These systems establish the architectural context for the
study; the evaluated contribution comprises the London data study, large-image
density and colour method, mesh-processing experiments, and ChordAtlas
integration path.

\section{Objectives}

Six objectives make the research question operational:

\begin{enumerate}
  \item construct and characterise a London multimodal case study comprising
  satellite imagery, panoramas, OSM context, and DSM data;
  \item adapt and evaluate large-image density fusion and per-vertex colour
  recovery while retaining the released Sat3DGen model;
  \item define the coordinate, tiling, and metadata contracts that connect
  geographic inputs, Sat3DGen density and mesh spaces, and ChordAtlas;
  \item implement mesh-space placement, cropping, overlap processing,
  DSM-informed height correction, and footprint-based building extraction;
  \item implement static and on-demand publication of building assets to
  ChordAtlas while preserving stable identity and vertex colour; and
  \item evaluate data integrity, geometric behaviour, topology, coordinate
  consistency, and cross-language publication, and relate the findings to
  Sat3DGen and procedural city-generation methods.
\end{enumerate}

\section{Contributions}

The first contribution is a London data-preparation and integrity study that
connects cross-view imagery with OSM and DSM context.  It quantifies geographic
support, channel availability, duplicate imagery, split relationships, and
candidate-label geometry.  These measurements show which properties must be
checked before a geographically assembled collection is used for learning or
evaluation.

The second contribution is a rectangular large-image inference method.  It
places overlapping density predictions at sub-voxel coordinates, blends them
with smooth spatial weights, extracts one global surface, and recovers
per-vertex RGB in a memory-bounded second pass.  A matched experiment measures
the reduction in window-boundary discontinuity relative to hard-crop fusion.

The third contribution is a mesh-space integration method for placing and
combining independently generated tiles.  The method introduces explicit
coordinate conversion, footprint association, DSM-informed vertical
processing, and per-building extraction.  Controlled geometry cases and
topological measurements evaluate the assumptions made by bottom removal,
stitching, and height correction.

The fourth contribution is a ChordAtlas bridge that publishes selected geometry
as modular building assets.  Stable request identities, manifests, validation
gates, staged publication, vertex-colour transfer, and Python--Java conversion
make the data boundary inspectable and repeatable.  The evaluation identifies
the conditions under which this bridge operates reliably and the geometry and
transaction properties that most strongly affect regional extension.

Sat3DGen, ChordAtlas, BigSUR, FrankenGAN, CityEngine, and Ubiq are pre-existing
systems.  The project-specific work reported here comprises the London data
workflow, large-image density and colour extension, geospatial mesh-processing
method, ChordAtlas integration, and their evaluation.  External software,
datasets, and published methods are identified and cited where they are used.

\section{Dissertation structure}

Chapter~\ref{chap:background} reviews procedural urban modelling,
satellite-conditioned reconstruction, semantic city representations,
geospatial contracts, and related systems.  Chapter~\ref{chap:methodology}
defines the data, large-image inference, coordinate, mesh-processing,
publication, and evaluation methods.  Chapter~\ref{chap:results} reports the
data-integrity, density-fusion, colour, geometry, coordinate, and ChordAtlas
integration results.  Chapter~\ref{chap:discussion}
interprets the findings, compares the approach with Sat3DGen and CityEngine,
and evaluates validity and future directions.  Chapter~\ref{chap:conclusion}
summarises the research outcome and contributions.  The appendix records
implementation provenance and detailed component-test results.
